\emph{The monotone zone, and finiteness of $S$.} By Lemma \claimref{sheaf.determinant-identity-blow-up-expansion}(2), $s = 1/\lambda + b$ with $b$ smooth up to the wall, locally uniformly; so $s' = -\lambda'/\lambda^2 + b'$, and since $\lambda \to 0$ while $\lambda'$ is bounded away from zero near $\Sig$ (T3, \claimref{sheaf.tameness}) and $b'$ is bounded (compactness of $\Sig$ makes all these bounds uniform along the wall), there is a one-sided collar neighbourhood on which $|\lambda'|/\lambda^2 > |b'|$: the monotone zone. Its far-side complement within the collar has compact closure disjoint from $\Sig$, on which the smooth function $s$ is bounded: $S < \infty$.

(1) Outside the collar $F = 1$; on the far collar side outside the monotone zone, $|2\mu\nu\beta^2 s| \leq 2|\mu\nu| S < 1$, so $F > 0$: no zeros in either region, whatever $\beta$ does there. All far-side zeros therefore lie in the monotone zone. Approaching the wall $F \to -\infty$ (Lemma \claimref{sheaf.determinant-identity-blow-up-expansion}(2)) and at the zone's outer edge $F > 0$, so each transversal carries a zero; where $\beta \equiv 1$ (shrink the zone into the inner collar if necessary), $F = 1 + 2\mu\nu s$ is strictly monotone along each transversal with a single simple zero, and $dF = 2\mu\nu\, ds$ has nonvanishing transversal component $2\mu\nu s' \neq 0$: $\{F = 0\}$ is a regular level set, a smooth closed hypersurface $\Sig'$, and $\det h = \det G \cdot (-F)$ has a simple zero across it ($\det G \neq 0$ there).

(2) At a simple zero of $\det h = \prod_i \mu_i$ (eigenvalues with respect to the auxiliary metric), exactly one eigenvalue $\mu_1$ vanishes and to first order: two vanishing eigenvalues would force a second-order zero. For the smoothness of the radical line, no spectral theory is needed: the adjugate $\operatorname{adj}(h)$, polynomial in the entries of $h$, is smooth; at a point of corank one it has rank exactly one (some $(d{-}1)$-minor is nonzero, and $h \cdot \operatorname{adj}(h) = \det(h)\,1 = 0$ puts every column of $\operatorname{adj}(h)$ in $\ker h$, which is a line). A nonvanishing column of $\operatorname{adj}(h)$ is therefore a smooth local spanning section of the radical line along $\Sig'$.

(3) Jacobi's formula, $\partial_\tau \det h = \operatorname{tr}\bigl(\operatorname{adj}(h)\, \partial_\tau h\bigr)$ (differentiate the determinant column by column and expand each along its differentiated column), evaluated at a point of $\Sig'$: in an auxiliary-orthonormal eigenbasis there, $\operatorname{adj}(h) = \bigl(\prod_{i \geq 2} \mu_i\bigr)\, \tilde k\, \tilde k^{\mathsf T}$ with $\tilde k$ the unit kernel vector, so
\[
\partial_\tau \det h \;=\; \Bigl(\prod_{i \geq 2} \mu_i\Bigr)\, (\partial_\tau h)(\tilde k, \tilde k) .
\]
The left side is nonzero (the zero of $\det h$ is simple, by (1)) and the nonvanishing eigenvalues $\mu_i$, $i \geq 2$, are nonzero; hence $(\partial_\tau h)(\tilde k, \tilde k) \neq 0$. This is exactly the nonvanishing of the datum $\delta'$ of Lemma \claimref{sheaf.invariant-form-t3} for $(T, h)$ along $\Sig'$: for any extension of $\tilde k$, $d\,h(\tilde k, \tilde k)|_{\Sig'} = (\partial_\tau h)(\tilde k, \tilde k)\, d\tau$, the cross term vanishing because $\tilde k$ is radical on $\Sig'$, and the lemma's proof uses only the symmetric form and the radical, not the cotangent origin of the bundle. So (T3) holds, (T1) and (T2) being parts (1) and (2).
